
\documentclass{article}
\usepackage{colortbl}
\usepackage{makecell}
\usepackage{multirow}
\usepackage{supertabular}

\begin{document}

\newcounter{utterance}

\centering \large Interaction Transcript for game `cladder', experiment `full\_v1.5\_default', episode 188 with refineonly{-}Qwen3.5{-}9B.
\vspace{24pt}

{ \footnotesize  \setcounter{utterance}{1}
\setlength{\tabcolsep}{0pt}
\begin{supertabular}{c@{$\;$}|p{.15\linewidth}@{}p{.15\linewidth}p{.15\linewidth}p{.15\linewidth}p{.15\linewidth}p{.15\linewidth}}
    \# & \multicolumn{2}{c}{Player} && \multicolumn{2}{c}{Game Master} \\
    \hline

    \theutterance \stepcounter{utterance}  
    & & & \multicolumn{4}{p{0.6\linewidth}}{
        \cellcolor[rgb]{0.9,0.9,0.9}{
            \makecell[{{p{\linewidth}}}]{
                \texttt{\tiny{[P1$\langle$GM]}}
                \texttt{You are an expert in causal inference. The following question is not a typical commonsense query, but rather a meticulously designed question created by a professor specializing in causal inference, intended to assess the students' mastery of the course content.} \\
\texttt{Imagine a self{-}contained, hypothetical world with only the following conditions, and without any unmentioned factors or causal relationships: The captain has a direct effect on the private and the corporal. The corporal has a direct effect on prisoner. The private has a direct effect on prisoner. Method 1: We look directly at how the captain correlates with prisoner in general. Method 2: We look at this correlation case by case according to the private. To understand how the captain affects prisoner, is it more correct to use the Method 1 than Method 2?} \\
\texttt{Start your answer with 'Yes' or 'No', followed by additional reasoning or evidence to support your explanation.} \\
            }
        }
    }
     \\ \\

    \theutterance \stepcounter{utterance}  
    & & \multicolumn{4}{p{0.6\linewidth}}{
        \cellcolor[rgb]{0.95,0.95,0.95}{
            \makecell[{{p{\linewidth}}}]{
                \texttt{\tiny{[P1$\leftrightarrow$P]}}
                \texttt{{[}DRAFT{]}} \\
\texttt{No.} \\
\\ 
\texttt{In this hypothetical world, the causal structure forms a Directed Acyclic Graph (DAG) where the **Captain** is a common cause of both the **Private** and the **Corporal**, and both the Private and the Corporal are direct causes of the **Prisoner**. This creates a "fork" structure at the Captain level (\$Captain \textbackslash{}to Private \textbackslash{}to Prisoner\$ and \$Captain \textbackslash{}to Corporal \textbackslash{}to Prisoner\$), resulting in a confounding path: \$Captain \textbackslash{}to Private \textbackslash{}to Prisoner \textbackslash{}leftarrow Corporal \textbackslash{}leftarrow Captain\$.} \\
\\ 
\texttt{**Method 1** calculates the marginal correlation between the Captain and the Prisoner. Because the Captain influences the Prisoner through two distinct pathways (via the Private and via the Corporal), the total effect observed in Method 1 is the sum of these direct and indirect effects. In causal inference, when the goal is to understand "how the captain affects prisoner" (i.e., the total causal effect), the marginal association is the correct metric, provided there are no unmeasured confounders between the Captain and the Prisoner other than the specified paths. Since the prompt states this is a self{-}contained world with no unmentioned factors, the marginal correlation captures the full causal impact.} \\
\\ 
\texttt{**Method 2** attempts to analyze the correlation "case by case according to the private." This phrasing implies conditioning on the Private (stratifying by the Private's status). Conditioning on a variable that lies on a causal pathway from the treatment (Captain) to the outcome (Prisoner) constitutes **overadjustment bias** (or blocking part of the causal mechanism). By fixing the Private's state, you artificially sever the pathway \$Captain \textbackslash{}to Private \textbackslash{}to Prisoner\$, thereby underestimating or distorting the total effect of the Captain. Furthermore, if the Private is a collider in a specific sub{-}context or if the conditioning introduces selection bias, Method 2 could yield a biased estimate of the total effect.} \\
\\ 
\texttt{Therefore, to understand the total effect of the Captain on the Prisoner, Method 1 is correct, while Method 2 introduces bias by conditioning on an intermediate variable.} \\
            }
        }
    }
    & \\ \\

    \theutterance \stepcounter{utterance}  
    & \multicolumn{4}{p{0.6\linewidth}}{
        \cellcolor[rgb]{0.561,0.737,0.561}{
            \makecell[{{p{\linewidth}}}]{
                \texttt{\tiny{[P1$\rangle$GM]}}
                \texttt{No. In this hypothetical world, the causal structure forms a Directed Acyclic Graph (DAG) where the **Captain** is a common cause of both the **Private** and the **Corporal**, and both the Private and the Corporal are direct causes of the **Prisoner**. This creates two distinct causal pathways: \$Captain \textbackslash{}to Private \textbackslash{}to Prisoner\$ and \$Captain \textbackslash{}to Corporal \textbackslash{}to Prisoner\$.} \\
\\ 
\texttt{**Method 1** calculates the marginal correlation between the Captain and the Prisoner. In causal inference, the marginal association between a treatment and an outcome captures the **total causal effect**, which is the sum of all direct and indirect pathways. Since the prompt specifies a self{-}contained world with no unmentioned confounders, the marginal correlation correctly reflects the full impact of the Captain on the Prisoner.} \\
\\ 
\texttt{**Method 2** attempts to analyze the correlation "case by case according to the private," which implies conditioning on the Private. Because the Private lies on the causal pathway from the Captain to the Prisoner, conditioning on this variable constitutes **overadjustment bias** (blocking part of the causal mechanism). By fixing the Private's status, Method 2 artificially severs the pathway \$Captain \textbackslash{}to Private \textbackslash{}to Prisoner\$, thereby underestimating the total effect of the Captain.} \\
\\ 
\texttt{Therefore, to understand "how the captain affects prisoner" (the total effect), Method 1 is the correct approach, while Method 2 introduces bias by conditioning on an intermediate variable.} \\
            }
        }
    }
    & & \\ \\

    \theutterance \stepcounter{utterance}  
    & & & \multicolumn{2}{p{0.3\linewidth}}{
        \cellcolor[rgb]{0.9,0.9,0.9}{
            \makecell[{{p{\linewidth}}}]{
                \texttt{\tiny{[GM$|$GM]}}
                \texttt{No} \\
            }
        }
    }
    & & \\ \\

    \theutterance \stepcounter{utterance}  
    & & & \multicolumn{2}{p{0.3\linewidth}}{
        \cellcolor[rgb]{0.9,0.9,0.9}{
            \makecell[{{p{\linewidth}}}]{
                \texttt{\tiny{[GM$|$GM]}}
                \texttt{yes} \\
            }
        }
    }
    & & \\ \\

    \theutterance \stepcounter{utterance}  
    & & & \multicolumn{2}{p{0.3\linewidth}}{
        \cellcolor[rgb]{0.9,0.9,0.9}{
            \makecell[{{p{\linewidth}}}]{
                \texttt{\tiny{[GM$|$GM]}}
                \texttt{game\_result = LOSE} \\
            }
        }
    }
    & & \\ \\

\end{supertabular}
}

\end{document}
